% =====================================================================
%  Part: the simple script, snippet by snippet.
%  Included from study_guide.tex.
% =====================================================================

\part{The simple script, snippet by snippet}
\label{part:simple}

This part walks through \texttt{fisherfaces\_simple.m} from the first line to
the last. Every block appears three times over:

\begin{enumerate}[nosep]
\item \textbf{the code}, copied verbatim from the file;
\item \textbf{the maths}, with the formula the code is computing;
\item \textbf{the syntax}, explaining every MATLAB construct on those lines.
\end{enumerate}

Nothing is assumed about MATLAB. Where a construct appears for the second time
it is not re-explained, but Section~\ref{sec:syntaxtable} collects all of them
in one table with a pointer back to where each is introduced.

%% =====================================================================
\section{Setting up}

\subsection{Clearing the workspace and fixing the random seed}

\begin{lstlisting}
clear; clc;
rng(0);
\end{lstlisting}

\paragraph{The maths.} Only the third line matters mathematically. The toy
problem in Part~1 draws $1000$ random points, so without a fixed seed every run
would print different numbers and no figure in this guide could be quoted.
Setting the seed makes the whole script a \emph{deterministic function} of the
data on disk.

\paragraph{The syntax.}
\begin{description}[leftmargin=*,style=nextline,nosep]
\item[\texttt{clear}] Deletes all variables from the workspace. Without it, a
  variable left over from a previous run could silently be picked up.
\item[\texttt{clc}] Clears the text in the command window. Purely cosmetic.
\item[\texttt{;}] A semicolon at the end of a line \emph{suppresses printing}
  of the result. Without it MATLAB echoes the value of every statement. This is
  the single most common surprise for people new to MATLAB.
\item[\texttt{rng(0)}] Resets the random number generator to a known state.
  Any fixed number works; $0$ is conventional.
\end{description}

\subsection{Finding the data}

\begin{lstlisting}
thisDir  = fileparts(mfilename('fullpath'));
if isempty(thisDir)
    thisDir = pwd;
end
orlRoot  = fullfile(thisDir, '..', 'ORL', filesep);
yaleRoot = fullfile(thisDir, '..', 'CroppedYale', filesep);
figDir   = fullfile(thisDir, 'figs_simple', filesep);
if exist(figDir, 'dir') == 0
    mkdir(figDir);
end
\end{lstlisting}

\paragraph{The maths.} None. This is bookkeeping, but it earns its place: an
earlier version used the relative path \texttt{'../ORL/'}, which only worked if
you happened to be sitting in the project folder. Run from anywhere else it
failed \emph{silently} --- see the gotcha below.

\paragraph{The syntax.}
\begin{description}[leftmargin=*,style=nextline,nosep]
\item[\texttt{mfilename('fullpath')}] Returns the full path of the script that
  is currently running, without the \texttt{.m} extension.
\item[\texttt{fileparts(p)}] Splits a path into directory, name and extension.
  Called with one output it returns just the directory.
\item[\texttt{isempty(x)}] True when \texttt{x} has zero elements. Here it
  catches the case of code pasted straight into the command window, where there
  is no file and \texttt{mfilename} returns \texttt{''}.
\item[\texttt{pwd}] ``Print working directory'' --- the current folder.
\item[\texttt{fullfile(a,b,c)}] Joins path pieces with the right separator for
  the operating system. Safer than gluing strings with \texttt{'/'}.
\item[\texttt{filesep}] The path separator character, \texttt{'/'} here.
\item[\texttt{exist(name,'dir')}] Returns $7$ if a folder of that name exists
  and $0$ if not. The \texttt{'dir'} argument restricts the question to folders.
\item[\texttt{mkdir}] Creates a folder.
\item[\texttt{if \ldots end}] MATLAB's \texttt{if} has no braces and no
  \texttt{then}; it is closed by \texttt{end}.
\end{description}

\begin{gotcha}[\texttt{dir} on a missing folder is not an error]
\texttt{dir('/nonexistent/*.pgm')} returns an \emph{empty list}, not an error.
So a wrong path does not stop the script: every loop simply runs zero times,
and the program carries on and ``classifies'' nothing. When this happened to
the big script it died three sections later inside the PCA with
\texttt{Index in position 2 exceeds array bounds} --- an error about the
wrong thing, in the wrong place. Both scripts now check the folders up front
and the loader raises an error if a subject folder yields no images.
\end{gotcha}

%% =====================================================================
\section{Part 1: the toy problem}

The point of this part is to make the difference between PCA and Fisher
visible in two dimensions, where you can check the answer by eye.

\subsection{Building two classes on purpose}

\begin{lstlisting}
Sigma = [16 0; 0 0.25];
mu1   = [ 0; -1];
mu2   = [ 0;  1];
nPer  = 500;

R  = chol(Sigma);
X1 = mu1 + R' * randn(2, nPer);
X2 = mu2 + R' * randn(2, nPer);
\end{lstlisting}

\paragraph{The maths.} We draw $n=500$ points from each of two Gaussians that
share a covariance matrix:
\[
  \mathcal{N}(\mu_1, \Sigma), \qquad \mathcal{N}(\mu_2, \Sigma), \qquad
  \Sigma = \begin{pmatrix} 16 & 0 \\ 0 & 0.25\end{pmatrix}, \quad
  \mu_1 = \begin{pmatrix}0\\-1\end{pmatrix}, \quad
  \mu_2 = \begin{pmatrix}0\\ 1\end{pmatrix}.
\]
The numbers are chosen to set a trap for PCA. Along $x$ the standard deviation
is $\sqrt{16}=4$; along $y$ it is $\sqrt{0.25}=0.5$. So the \emph{spread} is
$8$ times larger along $x$. But the two means differ \emph{only} in $y$. In
words: all of the variance is in the useless direction and all of the class
information is in the ignored one.

How the generation works: if $z \sim \mathcal{N}(0, I)$ then for any matrix $A$,
$Az$ has covariance
\[
  \operatorname{Cov}(Az) = A\,\operatorname{Cov}(z)\,A^{\!\top} = A I A^{\!\top} = AA^{\!\top}.
\]
We want $AA^{\!\top} = \Sigma$. Cholesky gives $\Sigma = R^{\!\top}R$ with $R$
upper triangular, so taking $A = R^{\!\top}$ gives
$AA^{\!\top} = R^{\!\top}R = \Sigma$, exactly what we need. Hence
\texttt{R' * randn(2,nPer)}, and adding $\mu$ shifts the cloud into place.

\paragraph{The syntax.}
\begin{description}[leftmargin=*,style=nextline,nosep]
\item[\texttt{[16 0; 0 0.25]}] A matrix literal. Spaces (or commas) separate
  entries in a row; a semicolon starts a new row. So this is $2\times2$.
\item[\texttt{[0; -1]}] A \emph{column} vector, because the semicolons stack
  the entries vertically. \texttt{[0 -1]} would be a row vector. In this whole
  project vectors are columns.
\item[\texttt{chol(A)}] Cholesky factorisation. Returns upper triangular
  \texttt{R} with \texttt{R'*R == A}. Requires \texttt{A} symmetric positive
  definite.
\item[\texttt{'} (apostrophe)] Transpose. \texttt{R'} is $R^{\!\top}$. (Strictly
  it is the conjugate transpose, but everything here is real.)
\item[\texttt{randn(2,n)}] A $2\times n$ matrix of standard normal random
  numbers.
\item[\texttt{mu1 + R'*randn(2,nPer)}] \textbf{Implicit expansion}: \texttt{mu1}
  is $2\times1$ and the other term is $2\times500$. MATLAB automatically
  repeats the column $500$ times to make the shapes match. Older code had to
  write \texttt{repmat(mu1,1,500)} for this.
\item[\texttt{*}] Matrix multiplication. Note \texttt{.*} (with a dot) would be
  element-by-element instead --- a different operation.
\end{description}

\subsection{The two scatter matrices}

\begin{lstlisting}
m1 = mean(X1, 2);
m2 = mean(X2, 2);
mAll = (m1 + m2) / 2;

D1 = X1 - m1;
D2 = X2 - m2;
Sw = D1 * D1' + D2 * D2';

Sb = nPer * (m1 - mAll) * (m1 - mAll)' + ...
     nPer * (m2 - mAll) * (m2 - mAll)';
\end{lstlisting}

\paragraph{The maths.} These are the definitions, written out for $c=2$
classes. The \emph{within-class scatter} measures how spread out each class is
around its own mean:
\[
  \Sw \;=\; \sum_{j=1}^{c} \sum_{x \in X_j} (x - \mu_j)(x - \mu_j)^{\!\top}.
\]
Writing $D_j$ for the matrix whose columns are the centred points
$x - \mu_j$, the inner sum is exactly $D_jD_j^{\!\top}$, which is why two
matrix products replace a double loop.

The \emph{between-class scatter} measures how far the class means are from the
overall mean, each weighted by how many points that class has:
\[
  \Sb \;=\; \sum_{j=1}^{c} n_j\,(\mu_j - \bar{\mu})(\mu_j - \bar{\mu})^{\!\top}.
\]
Each term $vv^{\!\top}$ is an \emph{outer product}: a column times a row,
giving a $2\times2$ matrix of rank $1$. Summing $c=2$ of them, and remembering
that $\mu_1 - \bar\mu = -(\mu_2 - \bar\mu)$, gives a matrix of rank $1$ here ---
an instance of the general bound $\rank(\Sb) \le c-1$.

\paragraph{The syntax.}
\begin{description}[leftmargin=*,style=nextline,nosep]
\item[\texttt{mean(X,2)}] The \texttt{2} is the \emph{dimension to average
  along}. \texttt{mean(X,2)} averages across columns and returns one column
  (the mean of the points); \texttt{mean(X,1)} would average down each column
  instead. Getting this backwards is a classic bug, and here it would average
  the two coordinates of each point together, which is meaningless.
\item[\texttt{X1 - m1}] Implicit expansion again: subtracts the $2\times1$ mean
  from all $500$ columns.
\item[\texttt{D1 * D1'}] $(2\times500)(500\times2) = 2\times2$. Note the order:
  \texttt{D1' * D1} would be $500\times500$ and mean something completely
  different.
\item[\texttt{v * v'}] Column times row $=$ outer product, a matrix.
  \texttt{v' * v} is row times column $=$ a single number, the inner product.
  One apostrophe decides which.
\item[\texttt{\ldots} (three dots)] Line continuation: the statement carries on
  to the next line.
\end{description}

\subsection{Fisher's direction}

\begin{lstlisting}
wLDA = Sw \ (m1 - m2);
wLDA = wLDA / norm(wLDA);
\end{lstlisting}

\paragraph{The maths.} We want the direction $w$ maximising the ratio of
projected between-class scatter to projected within-class scatter,
\[
  J(w) \;=\; \frac{w^{\!\top}\Sb\,w}{w^{\!\top}\Sw\,w}.
\]
For two classes this has a closed form. Write $\delta = \mu_1 - \mu_2$. With
equal class sizes $n$, a short calculation gives
$\Sb = \tfrac{n}{2}\,\delta\delta^{\!\top}$, so the numerator is
\[
  w^{\!\top}\Sb\,w \;=\; \tfrac{n}{2}\,w^{\!\top}\delta\delta^{\!\top}w
  \;=\; \tfrac{n}{2}\,(\delta^{\!\top}w)^2 ,
\]
using that $\delta^{\!\top}w$ is a scalar and so equals its own transpose.
Maximising $J$ therefore means maximising $(\delta^{\!\top}w)^2$ while holding
$w^{\!\top}\Sw w$ fixed. Setting the derivative of
$(\delta^{\!\top}w)^2 - \lambda\,w^{\!\top}\Sw w$ to zero, and using
$\partial(w^{\!\top}Aw)/\partial w = 2Aw$ for symmetric $A$, gives
\[
  2(\delta^{\!\top}w)\,\delta \;=\; 2\lambda \Sw w
  \qquad\Longrightarrow\qquad
  w \;=\; \frac{\delta^{\!\top}w}{\lambda}\,\Sw^{-1}\delta
  \;\propto\; \Sw^{-1}(\mu_1 - \mu_2).
\]
The leading factor is just a number, and $J(w)$ is unchanged when $w$ is
scaled --- both numerator and denominator are quadratic in $w$, so a factor
$t$ contributes $t^2$ to each and cancels. That is why we are free to
normalise $w$ to unit length on the second line, purely for readability.

Here it predicts
\[
  \Sigma^{-1}\delta \;=\; \mathrm{diag}\!\left(\tfrac{1}{16},\,4\right)
  \begin{pmatrix}0\\-2\end{pmatrix} \;=\; \begin{pmatrix}0\\-8\end{pmatrix},
\]
which is the $y$-axis. The code prints $(0.0045,\,-1.0000)$, confirming it.

\paragraph{The syntax.}
\begin{description}[leftmargin=*,style=nextline,nosep]
\item[\texttt{A \textbackslash\ b}] \textbf{Backslash, the single most
  important operator here.} It \emph{solves} the linear system $Ax = b$ and
  returns $x$. It does \textbf{not} form $A^{-1}$. Mathematically
  \texttt{A\textbackslash b} equals $A^{-1}b$, but it is computed by
  factorisation, which is both faster and more accurate. Writing
  \texttt{inv(Sw)*(m1-m2)} would give the same answer here and is considered
  bad practice.
\item[\texttt{norm(v)}] The Euclidean length $\|v\| = \sqrt{\sum_i v_i^2}$ for
  a vector.
\item[\texttt{/}] Ordinary division when the right operand is a scalar. (For
  matrices it is a \emph{right} solve --- see Section~\ref{sec:fisherfn}.)
\end{description}

\subsection{PCA's direction, for contrast}

\begin{lstlisting}
Xall = [X1, X2];
Dall = Xall - mean(Xall, 2);
St   = Dall * Dall';
[Vt, Dt] = eig(St);
[~, iBig] = max(diag(Dt));
wPCA = Vt(:, iBig);
\end{lstlisting}

\paragraph{The maths.} PCA ignores the labels entirely. It pools all $1000$
points, forms the \emph{total} scatter
\[
  \St \;=\; \sum_{i=1}^{N} (x_i - \bar{x})(x_i - \bar{x})^{\!\top},
\]
and takes the eigenvector of $\St$ with the largest eigenvalue, which is the
direction of greatest variance. Because the $x$-spread dominates, that is the
$x$-axis --- and the code prints $(-0.9998,\,-0.0187)$, which is the $x$-axis up
to an irrelevant sign.

This is the whole argument of the paper in miniature: the direction of largest
variance and the direction of best separation can be \emph{orthogonal}.

\paragraph{The syntax.}
\begin{description}[leftmargin=*,style=nextline,nosep]
\item[\texttt{[X1, X2]}] Horizontal concatenation: glues the two $2\times500$
  matrices into one $2\times1000$. A semicolon instead of a comma would stack
  them vertically.
\item[\texttt{[Vt, Dt] = eig(St)}] Two-output \texttt{eig}: \texttt{Vt} holds
  eigenvectors as \emph{columns}, and \texttt{Dt} is a \emph{diagonal matrix}
  of eigenvalues, not a vector. One-output \texttt{eig(St)} returns just a
  vector of eigenvalues.
\item[\texttt{diag(Dt)}] Pulls the diagonal out of a matrix as a column vector.
  Confusingly, \texttt{diag} applied to a \emph{vector} does the opposite and
  builds a diagonal matrix.
\item[\texttt{[\textasciitilde, iBig] = max(\ldots)}] \texttt{max} can return
  two things: the maximum value and \emph{where} it occurred. The tilde
  \texttt{\textasciitilde} means ``I do not want this output'', here the value.
  We only want the index.
\item[\texttt{Vt(:, iBig)}] Colon in the row slot means ``all rows'', so this
  is column \texttt{iBig} in full.
\end{description}

\begin{gotcha}[\texttt{eig} sorts ascending]
MATLAB's \texttt{eig} returns eigenvalues smallest-first. Every place in this
project that wants the \emph{largest} must sort or take a maximum explicitly.
Forgetting this does not cause an error --- it quietly returns the least
useful directions, and the recognition rate collapses to near chance.
\end{gotcha}

\subsection{Scoring the two directions}

\begin{lstlisting}
Jlda = (wLDA' * Sb * wLDA) / (wLDA' * Sw * wLDA);
Jpca = (wPCA' * Sb * wPCA) / (wPCA' * Sw * wPCA);
\end{lstlisting}

\paragraph{The maths.} This evaluates $J(w) = (w^{\!\top}\Sb w)/(w^{\!\top}\Sw w)$
for each candidate. Each of \texttt{w'*S*w} is a
$(1\times2)(2\times2)(2\times1) = 1\times1$ scalar. Measured:
$J(w_{\text{Fisher}}) = 4.24$ against $J(w_{\text{PCA}}) = 0.0021$, a factor of
about $2000$.

\paragraph{The syntax.} Nothing new: transpose, matrix product, scalar divide.
Note that MATLAB evaluates \texttt{w'*Sb*w} left to right, so it computes a
$1\times2$ row first and then a scalar --- never forming anything large.

\subsection{Turning $J$ into an error rate}

\begin{lstlisting}
errLDA = toyError(X1, X2, wLDA);
errPCA = toyError(X1, X2, wPCA);
\end{lstlisting}

with, at the bottom of the file,

\begin{lstlisting}
function err = toyError(X1, X2, w)
p1 = w' * X1;
p2 = w' * X2;
c1 = mean(p1);
c2 = mean(p2);

wrong = 0;
for i = 1:numel(p1)
    if abs(p1(i) - c1) > abs(p1(i) - c2)
        wrong = wrong + 1;
    end
end
for i = 1:numel(p2)
    if abs(p2(i) - c2) > abs(p2(i) - c1)
        wrong = wrong + 1;
    end
end
err = 100 * wrong / (numel(p1) + numel(p2));
end
\end{lstlisting}

\paragraph{The maths.} $J$ is an abstract ratio; an error rate is not. We
project every point onto the candidate direction,
\[
  p_i \;=\; w^{\!\top}x_i \qquad (\text{a single number per point}),
\]
compute the projected class means $c_1, c_2$, and assign each point to the
nearer one. A point is counted wrong when
$|p_i - c_{\text{own}}| > |p_i - c_{\text{other}}|$. Measured: $2.00\%$ along
Fisher's direction, $48.30\%$ along PCA's --- and $50\%$ is what a coin would
achieve. PCA has destroyed the problem.

\paragraph{The syntax.}
\begin{description}[leftmargin=*,style=nextline,nosep]
\item[\texttt{w' * X1}] $(1\times2)(2\times500) = 1\times500$: a row of
  projections, computed for all points in one multiplication.
\item[\texttt{function err = toyError(X1, X2, w)}] Declares a function: outputs
  on the left of \texttt{=}, inputs in brackets. Inside, \texttt{err} is an
  ordinary variable; whatever it holds at the end is returned.
\item[\texttt{for i = 1:numel(p1)}] \texttt{1:n} builds the vector
  $1,2,\ldots,n$ and the loop walks it. \texttt{numel} is the number of
  elements.
\item[\texttt{p1(i)}] Indexing is \textbf{1-based}: the first element is
  \texttt{p1(1)}, not \texttt{p1(0)}.
\item[\texttt{abs}] Absolute value.
\item[\texttt{end}] Closes the \texttt{for}, the \texttt{if}, \emph{and} the
  function. MATLAB uses the same keyword for all of them.
\end{description}

\subsection{The theoretical floor}

\begin{lstlisting}
Delta = sqrt((mu2 - mu1)' * (Sigma \ (mu2 - mu1)));
bayes = 100 * 0.5 * erfc((Delta / 2) / sqrt(2));
\end{lstlisting}

\paragraph{The maths.} For two Gaussians with a shared covariance and equal
priors, no classifier of any kind can do better than
\[
  P_{\text{error}} \;=\; \Phi\!\left(-\tfrac{\Delta}{2}\right),
  \qquad
  \Delta \;=\; \sqrt{(\mu_2-\mu_1)^{\!\top}\Sigma^{-1}(\mu_2-\mu_1)} ,
\]
where $\Delta$ is the \emph{Mahalanobis} distance between the means and $\Phi$
is the standard normal CDF. Here
\[
  \Delta = \sqrt{(0,2)\begin{pmatrix}1/16&0\\0&4\end{pmatrix}\begin{pmatrix}0\\2\end{pmatrix}}
  = \sqrt{16} = 4,
\]
so the floor is $\Phi(-2) = 2.28\%$. Our Fisher direction scored $2.00\%$.

\begin{gotcha}[Measuring \emph{below} the theoretical floor]
$2.00\% < 2.28\%$ looks like a contradiction and is not. The $2.28\%$ is an
expectation over the whole distribution; we measured on a finite sample of
$1000$ points, where $20$ errors instead of the expected $22.8$ is well inside
normal fluctuation. Quote this as sampling noise, not as beating the Bayes
bound --- nothing can beat the Bayes bound.
\end{gotcha}

\paragraph{The syntax.}
\begin{description}[leftmargin=*,style=nextline,nosep]
\item[\texttt{Sigma \textbackslash\ v}] Backslash again: solves
  $\Sigma x = v$, giving $\Sigma^{-1}v$ without inverting.
\item[\texttt{erfc(z)}] The complementary error function. MATLAB has no
  \texttt{normcdf} without the Statistics Toolbox, which is not installed here,
  so we use the identity $\Phi(-t) = \tfrac12\operatorname{erfc}(t/\sqrt{2})$.
\end{description}

%% =====================================================================
\section{Part 2: real faces}

\subsection{Listing the subject folders}

\begin{lstlisting}
orlFolders = cell(1, 32);
for p = 1:32
    orlFolders{p} = [orlRoot 's' num2str(p) '/'];
end

yaleFolders = cell(1, 38);
for p = 1:38
    yaleFolders{p} = [yaleRoot 's' num2str(p) '/'];
end
\end{lstlisting}

\paragraph{The maths.} None, but one decision matters: \textbf{the class label
is the loop index \texttt{p}}, never the folder's name. That is what makes the
folder names arbitrary. The Yale folders ship as
\texttt{yaleB01}\ldots\texttt{yaleB39} with \texttt{yaleB14} missing, which is
not a countable range; \texttt{rename\_yale.sh} renames them to
\texttt{s1}\ldots\texttt{s38} so both databases load with the same loop.
Renaming cannot change any result, because recognition depends only on
\emph{which images share a label}, and the alphabetical order of the old names
maps onto $1\ldots38$ in the same order.

\paragraph{The syntax.}
\begin{description}[leftmargin=*,style=nextline,nosep]
\item[\texttt{cell(1,32)}] Creates a $1\times32$ \emph{cell array}. A normal
  array holds numbers of one type; a cell array holds anything, here strings of
  different lengths.
\item[\texttt{orlFolders\{p\}}] \textbf{Braces} index a cell array and give you
  \emph{the contents}. Round brackets \texttt{orlFolders(p)} would give you a
  $1\times1$ cell still wrapped up. This distinction causes a lot of confusion;
  the rule is: braces to get the thing out.
\item[\texttt{['abc' 'def']}] Square brackets concatenate strings, because a
  character string is just a row vector of characters. So
  \texttt{[orlRoot 's' num2str(p) '/']} glues four pieces into one path.
\item[\texttt{num2str(p)}] Converts the number $7$ into the characters
  \texttt{'7'}, so it can be glued into a path.
\end{description}

\begin{gotcha}[Why ORL is not globbed]
It would be natural to write \texttt{dir([orlRoot 's*'])}. Do not: this copy of
ORL contains a stray empty directory called \texttt{sed}, which matches
\texttt{s*} and would appear as a $41$st subject. Building the path from the
index never sees it.
\end{gotcha}

\subsection{Loading the images}
\label{sec:loadfn}

\begin{lstlisting}
nTr = 0;
nTe = 0;
for p = 1:nF
    f = dir([folders{p} '*.pgm']);
    if isempty(f)
        error('No .pgm images found in %s', folders{p});
    end
    ntr = min(numTrain, numel(f));
    nTr = nTr + ntr;
    nTe = nTe + (numel(f) - ntr);
end

imTr = zeros(d, nTr);   ytr = zeros(1, nTr);
imTe = zeros(d, nTe);   yte = zeros(1, nTe);

a = 0;
b = 0;
for p = 1:nF
    files = dir([folders{p} '*.pgm']);
    for j = 1:numel(files)
        im = double(imread([folders{p} files(j).name]));
        if j <= numTrain
            a = a + 1;
            imTr(:, a) = im(:);
            ytr(a)     = p;
        else
            b = b + 1;
            imTe(:, b) = im(:);
            yte(b)     = p;
        end
    end
end
\end{lstlisting}

\paragraph{The maths.} The modelling decision here is \emph{vectorisation of an
image}. An image is a grid of $H \times W$ pixel intensities; we stack its
columns into one long vector
\[
  x \in \R^{d}, \qquad d = HW ,
\]
and from then on treat it as a point in $d$-dimensional space. For ORL
$d = 112\cdot92 = 10304$; for Yale $d = 192\cdot168 = 32256$. This throws away
the knowledge that pixel $(i,j)$ is next to pixel $(i,j{+}1)$ --- the methods
here never use spatial adjacency. It is legitimate because every image is
unrolled the \emph{same} way, so corresponding pixels land in corresponding
coordinates, and that is all the geometry we need.

The training set becomes a matrix with \textbf{one image per column},
\[
  \mathrm{imTr} = \begin{pmatrix} | & & | \\ x_1 & \cdots & x_N \\ | & & | \end{pmatrix}
  \in \R^{d\times N},
\]
with $\mathrm{ytr}(i)$ the identity of the person in column $i$.

\paragraph{Where the labels come from.} This is worth stating plainly, because
LDA is a supervised method and the question ``where did the labels come from?''
has a mundane answer: \emph{the folder each file was in}. We never look inside
an image to decide whose face it is.

\paragraph{The syntax.}
\begin{description}[leftmargin=*,style=nextline,nosep]
\item[\texttt{dir(pattern)}] Returns a \emph{struct array}, one entry per
  matching file, with fields \texttt{name}, \texttt{folder}, \texttt{bytes} and
  so on. \texttt{files(j).name} is the $j$-th file's name.
\item[\texttt{zeros(d,n)}] Allocates a $d\times n$ matrix of zeros. We do this
  \emph{before} the loop (``preallocation''): growing an array inside a loop
  makes MATLAB copy the whole thing every iteration, and for Yale that would be
  $1520$ copies of a $392$\,MB matrix.
\item[\texttt{imread(path)}] Reads an image file into a matrix.
\item[\texttt{double(\ldots)}] Converts from \texttt{uint8} (integers
  $0$--$255$) to double precision. \textbf{Essential}: in \texttt{uint8}
  arithmetic, $200 + 100$ saturates to $255$ and $5 - 10$ clips to $0$, so
  every mean and difference would be silently wrong.
\item[\texttt{im(:)}] \textbf{The colon on its own reshapes anything into a
  single column}, stacking \emph{column by column}. MATLAB is column-major, so
  a $2\times2$ matrix $\left(\begin{smallmatrix}1&3\\2&4\end{smallmatrix}\right)$ becomes
  $(1,2,3,4)^{\!\top}$. The inverse is \texttt{reshape(v,H,W)}, used in Part~3.
\item[\texttt{imTr(:, a) = \ldots}] Assigns into \emph{all rows} of column
  \texttt{a}.
\item[\texttt{error('\ldots \%s', x)}] Raises an error, with \texttt{\%s}
  replaced by \texttt{x}, stopping the script.
\item[\texttt{min(a,b)}] The smaller of two numbers --- here guarding against a
  subject with fewer images than \texttt{numTrain}.
\end{description}

\begin{gotcha}[\texttt{dir} sorts names as \emph{text}]
ORL's files are \texttt{1.pgm, 2.pgm, \ldots, 10.pgm}, and text order puts
\texttt{10} straight after \texttt{1}. So ``the first six'' is really
$\{1,10,2,3,4,5\}$. For Yale the names encode lighting as \texttt{A+020} and
\texttt{A-020}, and \texttt{'+'} sorts before \texttt{'-'}, so taking the first
$40$ files trains on \emph{positive} azimuths and tests on \emph{negative}
ones. That single accident roughly halves every Yale number in this guide ---
see Section~\ref{sec:splitbias}.
\end{gotcha}

\subsection{The correlation baseline}

\begin{lstlisting}
results(s, 1) = nnClassify(normCols(imTr), ytr, normCols(imTe), yte);
\end{lstlisting}

\begin{lstlisting}
function Xn = normCols(X)
Xn = X - mean(X, 1);
for i = 1:size(Xn, 2)
    nrm = norm(Xn(:, i));
    if nrm > 0
        Xn(:, i) = Xn(:, i) / nrm;
    end
end
end
\end{lstlisting}

\paragraph{The maths.} This is the ``do nothing clever'' baseline: nearest
neighbour on the raw pixels. Each image is first made zero-mean and unit-norm,
\[
  \tilde{x} \;=\; \frac{x - \operatorname{mean}(x)\mathbf{1}}{\|x - \operatorname{mean}(x)\mathbf{1}\|} .
\]
That turns Euclidean nearest neighbour into \emph{correlation} matching,
because for unit vectors
\[
  \|\tilde a - \tilde b\|^2
  = \tilde a^{\!\top}\tilde a - 2\tilde a^{\!\top}\tilde b + \tilde b^{\!\top}\tilde b
  = 2 - 2\,\tilde a^{\!\top}\tilde b ,
\]
so the smallest distance is attained by the largest $\tilde a^{\!\top}\tilde b$,
which is exactly the correlation coefficient. Subtracting the mean removes a
brightness offset and dividing by the norm removes a contrast scaling, so the
match is invariant to both.

\begin{observation}
Skipping the normalisation changes the answer and not in a consistent
direction: ORL reads $93.8\%$ instead of $91.4\%$, Yale $35.6\%$ instead of
$43.0\%$. It is not a cosmetic step.
\end{observation}

\paragraph{The syntax.}
\begin{description}[leftmargin=*,style=nextline,nosep]
\item[\texttt{mean(X,1)}] Here the \texttt{1} means average \emph{down} each
  column, i.e.\ the mean intensity of each image. Contrast with
  \texttt{mean(X,2)} earlier, which averaged across images.
\item[\texttt{size(X,2)}] The number of columns. \texttt{size(X,1)} is rows.
\end{description}

\subsection{PCA, and the \texorpdfstring{$L = X^{\!\top}X$}{L = X'X} trick}
\label{sec:pcafn}

\begin{lstlisting}
[~, N] = size(imgs);

xbar = mean(imgs, 2);
X    = imgs - xbar;

L = X' * X;
[Wg, D] = eig(L);

[eigvals, idx] = sort(diag(D), 'descend');
Wg = Wg(:, idx);

V = X * Wg;

scale = norm(V(:, 1));
for l = 1:size(V, 2)
    nrm = norm(V(:, l));
    if nrm > 1e-10 * scale
        V(:, l) = V(:, l) / nrm;
    else
        V(:, l) = 0;
    end
end

lambda = eigvals / N;
\end{lstlisting}

\paragraph{The maths.} PCA wants the top eigenvectors of the covariance
\[
  C \;=\; \frac{1}{N}XX^{\!\top} \in \R^{d\times d},
  \qquad X = \begin{pmatrix}x_1 - \bar x & \cdots & x_N - \bar x\end{pmatrix}.
\]
For Yale $C$ would be $32256\times32256$, which is $8.3$\,GB --- we cannot even
store it. But we do not need all of it.

\begin{theorem}[The Gram trick]
Let $L = X^{\!\top}X \in \R^{N\times N}$. If $Lw = \lambda w$ with $w \neq 0$
and $\lambda \neq 0$, then $Xw$ is an eigenvector of $XX^{\!\top}$ with the same
eigenvalue $\lambda$.
\end{theorem}
\begin{proof}
Multiply $X^{\!\top}Xw = \lambda w$ on the left by $X$:
\[
  X(X^{\!\top}Xw) \;=\; X(\lambda w)
  \qquad\Longrightarrow\qquad
  (XX^{\!\top})(Xw) \;=\; \lambda\,(Xw),
\]
where the left side was regrouped using associativity of matrix
multiplication. So $Xw$ satisfies the eigenvector equation for $XX^{\!\top}$.
It is non-zero because $\|Xw\|^2 = w^{\!\top}X^{\!\top}Xw = \lambda\|w\|^2 \neq 0$.
\end{proof}

So we solve a $192\times192$ problem instead of a $10304\times10304$ one, and
multiply by $X$ to get back. The same calculation shows
\[
  \|Xw\| \;=\; \sqrt{\lambda}\,\|w\| \;=\; \sqrt{\lambda} \quad (\text{for unit } w),
\]
which is \emph{not} $1$ --- hence the renormalisation loop. Without it
$V^{\!\top}V \neq I$ and the projection is no longer an orthogonal projection.

Finally, eigenvalues of $L = X^{\!\top}X$ are $N$ times those of
$C = XX^{\!\top}/N$, which is why \texttt{lambda = eigvals / N}.

\paragraph{Why centre first?} PCA finds directions of maximum \emph{variance},
which is spread about the centre. Without subtracting $\bar x$, the first
component would point at the mean face itself --- the direction in which the
raw vectors are longest --- describing what faces have in common rather than
how they differ.

\begin{gotcha}[The last direction is numerically empty]
After centring, the columns of $X$ sum to zero, so $\rank(X) \le N-1$ and the
smallest eigenvalue is zero up to round-off. Its ``eigenvector'' is noise with
a tiny norm, and dividing by that norm would amplify pure round-off into a
unit-length vector that looks like a legitimate eigenface. The guard
\texttt{nrm > 1e-10 * scale} is \emph{relative} to the largest column, so it
scales with the data; an absolute threshold failed here, because the measured
norm of the empty column was $1.95\times10^{-9}$, comfortably above a fixed
$10^{-10}$.
\end{gotcha}

\paragraph{The syntax.}
\begin{description}[leftmargin=*,style=nextline,nosep]
\item[\texttt{[\textasciitilde, N] = size(imgs)}] Two-output \texttt{size}
  returns rows and columns; we discard the rows with \texttt{\textasciitilde}.
\item[\texttt{sort(v,'descend')}] Sorts largest-first. With two outputs,
  \texttt{[vals, idx] = sort(\ldots)}, \texttt{idx} records where each sorted
  value came from.
\item[\texttt{Wg(:, idx)}] Reorders \emph{columns} using an index vector.
  Because \texttt{idx} came from sorting the eigenvalues, this puts the
  eigenvectors into the matching order. Doing one without the other is a
  nasty, silent bug.
\item[\texttt{1e-10}] Scientific notation for $10^{-10}$.
\end{description}

\subsection{Classifying with eigenfaces}

\begin{lstlisting}
ks   = [10 20 30 50 80 120 170];
ks   = ks(ks <= N - 1);
best = 0;  bestK = 0;
for k = ks
    Vk = V(:, 1:k);
    r  = nnClassify(Vk' * (imTr - xbar), ytr, Vk' * (imTe - xbar), yte);
    if r > best
        best  = r;
        bestK = k;
    end
end
\end{lstlisting}

\paragraph{The maths.} Keeping the top $k$ eigenfaces $V_k$, each image is
represented by its coefficient vector
\[
  a \;=\; V_k^{\!\top}(x - \bar{x}) \;\in\; \R^{k},
\]
and we classify in that $k$-dimensional space instead of in $\R^d$. The
compression is dramatic: $10304$ numbers become $50$.

\paragraph{The syntax.}
\begin{description}[leftmargin=*,style=nextline,nosep]
\item[\texttt{ks(ks <= N-1)}] \textbf{Logical indexing.} The inner expression
  \texttt{ks <= N-1} produces a vector of true/false, one per element; using it
  as an index keeps only the elements where it is true. This is MATLAB's
  idiomatic filter, replacing a loop with an \texttt{if} inside.
\item[\texttt{V(:, 1:k)}] A \emph{range} of columns: columns $1$ through $k$.
\item[\texttt{for k = ks}] A \texttt{for} loop can walk the elements of any
  vector, not just \texttt{1:n}.
\end{description}

\subsection{Fisherfaces}
\label{sec:fisherfn}

\begin{lstlisting}
Vm = V(:, 1:m);
Al = Vm' * (imgs - xbar);
abar = mean(Al, 2);

Sw = zeros(m, m);
Sb = zeros(m, m);
for j = 1:c
    Aj = Al(:, y == j);
    nj = size(Aj, 2);
    mj = mean(Aj, 2);
    Dj = Aj - mj;
    Sw = Sw + Dj * Dj';
    Sb = Sb + nj * (mj - abar) * (mj - abar)';
end

Sw = (Sw + Sw') / 2;
Sb = (Sb + Sb') / 2;
\end{lstlisting}

\paragraph{The maths.} Exactly the same two definitions as in the toy problem,
but now in the $m$-dimensional PCA space rather than in $\R^d$, and with a loop
over all $c$ classes instead of two written out. The reason we work in the
reduced space is the whole point of the method:

\begin{quote}
In pixel space $\Sw$ is $d\times d$ but $\rank(\Sw) \le N - c$. For Yale that
is $1482 \ll 32256$, so $\Sw$ is singular, $\Sw^{-1}$ does not exist, and
$J(w) = (w^{\!\top}\Sb w)/(w^{\!\top}\Sw w)$ is undefined for any $w$ in the
null space. Projecting onto $m \le N-c$ PCA directions discards exactly the
directions $\Sw$ has no information about; what remains has full rank $m$ and
is invertible.
\end{quote}

The rank bound itself is easy to see: class $j$ contributes $n_j$ vectors
$x - \mu_j$ that are forced to sum to zero, so it supplies at most $n_j - 1$
independent directions, and $\sum_j (n_j - 1) = N - c$.

\paragraph{The three rank results this code depends on.} These are proved in
Section~\ref{sec:rankthms}. They are worth knowing as a group, because between
them they decide the shape of everything below: how big $m$ may be, how many
Fisherfaces come out, and why ORL is the hard database.

\begin{center}
\small
\begin{tabular}{llll}
\toprule
 & statement & what it controls & measured \\
\midrule
Thm~\ref{thm:rankSb}  & $\rank(\Sb)\le c-1$ & how many directions carry
  & ORL: $31$, bound $31$ \\
 & & any class information & (tight) \\[3pt]
Cor~\ref{cor:ceiling} & at most $c-1$ Fisherfaces & the width of \texttt{Wopt}
  & ORL $31$, Yale $37$ \\[3pt]
Thm~\ref{thm:rankSw}  & $\rank(\Sw)\le N-c$ & the largest usable $m$
  & Yale: $1482$, bound $1482$ \\
\bottomrule
\end{tabular}
\end{center}

\begin{observation}[What each one looks like in the output]
\emph{Theorem~\ref{thm:rankSb} is tight, and the bound appears as a cliff rather
than a decay.} On ORL the generalised eigenvalues run
$\lambda_{1} = 1.49\times10^{7}$ down to $\lambda_{31} = 7.47$, and then
$\lambda_{32} = 2.1\times10^{-12}$ --- twelve orders of magnitude between two
consecutive eigenvalues, exactly at the predicted index. Check~(b) in the
script confirms $\rank(\Sb) = 31$ with $c = 32$.

\emph{Corollary~\ref{cor:ceiling} is why \texttt{Z(:, 1:(c-1))} is not a tuning
choice.} Keeping more columns cannot help: past $c-1$ the eigenvalues are zero
and their eigenvectors are an arbitrary basis of a null space. This also makes
the fair comparison against eigenfaces obvious --- Fisherfaces get exactly
$c-1$ features, so the script reports eigenfaces at $k = c-1$ as well as at
their best $k$.

\emph{Theorem~\ref{thm:rankSw} is the one that predicts the headline result.}
Read it as degrees of freedom: each person donates only $n_j - 1$ directions to
$\Sw$. ORL gives $6$ training images per person, so $\mathbf{5}$; Yale gives
$40$, so $\mathbf{39}$. $\Sw$ is our estimate of ``how one person's images
vary'', and on ORL it is built from five directions per person. That single
number is why Fisherfaces \emph{lose} on ORL ($90.6\%$ against eigenfaces'
$95.3\%$) and win on Yale ($67.2\%$ against $34.2\%$).
\end{observation}

\begin{gotcha}[A rank one short of the bound is a fact about your data]
$N-c$ is an \emph{upper} bound, and real data can fall short of it. On our
first Yale run the measured $\rank(\Sw)$ was $1481$, not $1482$. A deficiency
beyond the predicted one means some class contains an \emph{exact} linear
dependency, and the simplest cause is a repeated image. A duplicate search
found one pair in subject 12: a stray file
\texttt{yaleB12\_P00A+020E+10 - Copy.pgm}, byte-identical to the original. Two
identical columns stay identical after within-class centring, so they cost
exactly one unit of rank: $1482 - 1 = 1481$, which is what we measured. After
deleting it the rank meets the bound exactly.

The lesson generalises: a rank that disagrees with theory by a small amount is
usually telling you something true about the data, not about the algorithm.
\end{gotcha}

\paragraph{The syntax.}
\begin{description}[leftmargin=*,style=nextline,nosep]
\item[\texttt{Al(:, y == j)}] \textbf{Logical indexing again, now selecting
  columns.} \texttt{y == j} is a true/false row with one entry per image, true
  where that image belongs to person $j$; using it in the column slot keeps
  exactly those columns. This replaces a loop that would test each image.
  Note \texttt{==} is comparison; \texttt{=} is assignment.
\item[\texttt{(Sw + Sw')/2}] Forces exact symmetry. In exact arithmetic
  $D D^{\!\top}$ is already symmetric, but floating-point addition is not
  associative, so the computed matrix differs from its transpose by around
  $10^{-16}$. \texttt{chol} and symmetric \texttt{eig} read only \emph{one
  triangle}, so without this they would silently use half the matrix and we
  would not know which.
\end{description}

\subsubsection{Solving the generalised eigenproblem}

\begin{lstlisting}
R = chol(Sw);
M = (R') \ (Sb / R);
M = (M + M') / 2;

[Z, Dz] = eig(M);
[~, idx] = sort(diag(Dz), 'descend');
Z = Z(:, idx);

Wfld = R \ Z(:, 1:(c-1));
Wopt = Vm * Wfld;
\end{lstlisting}

\paragraph{The maths.} For more than two classes the maximiser of $J$ is no
longer a single closed-form vector. Instead the stationary points of
\[
  J(w) \;=\; \frac{w^{\!\top}\Sb\,w}{w^{\!\top}\Sw\,w}
\]
satisfy the \textbf{generalised eigenvalue problem}
\[
  \Sb\,w \;=\; \lambda\,\Sw\,w .
\]
We could multiply by $\Sw^{-1}$ and solve $\Sw^{-1}\Sb w = \lambda w$, but
$\Sw^{-1}\Sb$ is not symmetric --- its transpose is $\Sb\Sw^{-1}$ --- so we
would lose the guarantees of a symmetric solver. Instead, \emph{whiten}.

Since $\Sw$ is symmetric positive definite it factors as $\Sw = R^{\!\top}R$
with $R$ upper triangular and invertible. Substitute $w = R^{-1}z$:
\[
  \Sb R^{-1}z \;=\; \lambda\,\Sw R^{-1}z \;=\; \lambda\,R^{\!\top}R R^{-1} z
  \;=\; \lambda\,R^{\!\top}z .
\]
Multiply on the left by $R^{-\top} := (R^{\!\top})^{-1}$:
\[
  \boxed{\;\underbrace{R^{-\top}\Sb R^{-1}}_{=:\,M}\,z \;=\; \lambda z\;}
\]
and $M$ is symmetric, because
$M^{\!\top} = (R^{-\top}\Sb R^{-1})^{\!\top} = R^{-\top}\Sb^{\!\top}R^{-1} = M$
using $\Sb^{\!\top} = \Sb$. So we have converted a generalised problem into an
\emph{ordinary symmetric} one, where eigenvalues are guaranteed real and
eigenvectors orthogonal. Having found $z$, recover $w = R^{-1}z$.

The eigenvalues are untouched by the substitution, so each $\lambda_l$ is still
exactly the Fisher criterion value $J(w_l)$.

\paragraph{Why only $c-1$ columns.} $\Sb$ is a sum of $c$ rank-one terms whose
arguments $\mu_j - \bar\mu$ sum to zero after weighting, so $\rank(\Sb)\le c-1$.
Every eigenvalue past index $c-1$ is therefore zero, and its eigenvector is an
arbitrary basis vector of a null space --- noise. Measured on ORL, the drop is
not gradual but a cliff: $\lambda_{31} = 7.47$ and $\lambda_{32} = 2.1\times10^{-12}$.

\paragraph{The syntax.}
\begin{description}[leftmargin=*,style=nextline,nosep]
\item[\texttt{Sb / R}] \textbf{Forward slash is a \emph{right} solve}: it
  computes $\Sb R^{-1}$ by solving $YR = \Sb$ for $Y$. Compare backslash, which
  solves from the left. Together,
  \texttt{(R')\textbackslash(Sb/R)} computes $R^{-\top}\Sb R^{-1}$ with
  \emph{two triangular solves and no inverse anywhere}. Triangular systems are
  solved by simple substitution, so this is both fast and accurate.
\item[\texttt{Z(:, 1:(c-1))}] Keeps the first $c-1$ columns, i.e.\ those
  belonging to the largest eigenvalues after the sort.
\end{description}

\subsubsection{The normalisation that is not cosmetic}

\begin{lstlisting}
for l = 1:size(Wopt, 2)
    nrm = norm(Wopt(:, l));
    if nrm > 0
        Wopt(:, l) = Wopt(:, l) / nrm;
    end
end
\end{lstlisting}

\paragraph{The maths.} $J(w)$ is invariant under scaling, so mathematically the
subspace we have found is already correct and the lengths of its basis vectors
are arbitrary. But our \emph{classifier} is nearest neighbour with Euclidean
distance, and that is \textbf{not} scale invariant: in
\[
  \|a - b\|^2 \;=\; \sum_{l} (a_l - b_l)^2 ,
\]
a coordinate whose basis vector is long contributes more to every distance.
The columns returned by \texttt{R\textbackslash Z} satisfy
$W^{\!\top}\Sw W = I$, which forces their Euclidean norms to differ by orders
of magnitude --- whitening scales direction $l$ by $1/\sqrt{\lambda_l}$, and on
ORL the extreme ratio is about $8860$. Leaving them unnormalised lets a few
directions dominate every comparison.

\begin{observation}
Deleting these four lines costs ORL $90.6\% \to 20.3\%$ and Yale
$67.2\% \to 13.5\%$. The subspace, the theory and the eigenvalues are
identical in both cases; only the basis lengths differ.
\end{observation}

\subsection{The three checks}

\begin{lstlisting}
Al = V(:, 1:m)' * (imTr - xbar);
Dl = Al - mean(Al, 2);
St = Dl * Dl';
norm(St - (Sb + Sw), 'fro') / norm(St, 'fro')

rank(Sb)

Ge = Wopt' * Wopt;
offEuclid = norm(Ge - diag(diag(Ge)), 'fro') / norm(Ge, 'fro');
Wfld = V(:, 1:m)' * Wopt;
Gs   = Wfld' * Sw * Wfld;
offSw = norm(Gs - diag(diag(Gs)), 'fro') / norm(Gs, 'fro');
\end{lstlisting}

\noindent(In the file the first two quantities sit inside the \texttt{fprintf}
calls that print them; they are lifted out here so the three checks read as
three calculations. Every other listing in this part is copied verbatim.)

\paragraph{Check (a): $\St = \Sb + \Sw$.} A theorem, so it must hold to machine
precision. Expanding the total scatter and inserting $\pm\mu_j$:
\[
  \St = \sum_j \sum_{x\in X_j} \big((x - \mu_j) + (\mu_j - \bar\mu)\big)
                              \big((x - \mu_j) + (\mu_j - \bar\mu)\big)^{\!\top},
\]
the two cross terms vanish because $\sum_{x \in X_j}(x - \mu_j) = 0$ by the
definition of $\mu_j$, leaving $\Sw + \Sb$. Measured: $5.5\times10^{-16}$ on
ORL. A wrong class mean or a miscounted $n_j$ would break it immediately, so it
is a free correctness test on the whole scatter construction.

\paragraph{Check (b): $\rank(\Sb) = c-1$.} This is
Theorem~\ref{thm:rankSb} and Corollary~\ref{cor:ceiling} being tested at
runtime. Measured $31$ with $c = 32$, and we keep exactly $31$ Fisherfaces, so
the bound is tight here rather than merely an inequality. Were it to come out
\emph{below} $c-1$, that would not be a bug in the solver but a statement about
the data --- two class means coinciding, or a class with a single image.

\paragraph{Check (c): $\Sw$-orthogonality.} Fisherfaces are orthogonal in the
inner product $\langle u,v\rangle_{\Sw} = u^{\!\top}\Sw v$, \textbf{not} in the
ordinary one. The code measures the share of energy sitting off the diagonal of
each Gram matrix:
\[
  \frac{\|G - \operatorname{diag}(G)\|_F}{\|G\|_F},
  \qquad G = W^{\!\top}W \ \text{ or } \ W^{\!\top}\Sw W .
\]
Measured on ORL: $0.499$ for the Euclidean Gram --- \emph{half the energy is
off-diagonal} --- against $3.8\times10^{-13}$ for the $\Sw$ Gram. So if asked
``are Fisherfaces orthogonal?'', the answer is: in the $\Sw$ inner product yes,
in the usual one emphatically no.

\paragraph{The syntax.}
\begin{description}[leftmargin=*,style=nextline,nosep]
\item[\texttt{norm(A,'fro')}] The Frobenius norm,
  $\sqrt{\sum_{ij}A_{ij}^2}$ --- the natural ``size'' of a matrix. Plain
  \texttt{norm(A)} for a matrix would give the largest singular value instead.
\item[\texttt{diag(diag(G))}] The double \texttt{diag} is the idiom for
  ``keep only the diagonal''. The inner call extracts the diagonal as a vector,
  the outer rebuilds a matrix from it.
\item[\texttt{rank(A)}] Counts singular values above a tolerance. It is a
  \emph{numerical} judgement, not an exact one --- see the gotcha below.
\end{description}

\begin{gotcha}[\texttt{rank} and \texttt{chol} can disagree, and both be right]
On the Yale split, $\Sw$ in the reduced space has $m = 1482$ and
$\kappa(\Sw) = 6.0\times10^{12}$. \texttt{rank} reports $1481$ --- apparently
singular --- while \texttt{chol} succeeds, which requires positive
definiteness. Both are correct: \texttt{rank}'s default tolerance is
$\max(\text{size})\cdot\varepsilon\cdot\lambda_{\max}
= 1482 \cdot 2.22\times10^{-16}\lambda_{\max} = 3.3\times10^{-13}\lambda_{\max}$,
while the smallest eigenvalue sits at $1.7\times10^{-13}\lambda_{\max}$ --- below
\texttt{rank}'s threshold but above what Cholesky needs. This is what
``$m = N-c$ has zero slack'' looks like in practice.
\end{gotcha}

\subsection{Nearest neighbour}

\begin{lstlisting}
function rate = nnClassify(trainProj, trainLabels, testProj, testLabels)
correct = 0;
for i = 1:size(testProj, 2)
    distances = sum((trainProj - testProj(:, i)).^2, 1);
    [~, best] = min(distances);
    if trainLabels(best) == testLabels(i)
        correct = correct + 1;
    end
end
rate = 100 * correct / size(testProj, 2);
end
\end{lstlisting}

\paragraph{The maths.} The classifier assigns a test point the label of its
closest training point:
\[
  \hat{y}(x) \;=\; y_{i^*}, \qquad
  i^* \;=\; \arg\min_{i} \;\|a(x) - a_i\|^2 ,
\]
where $a(\cdot)$ is whichever projection is being tested. There is
\textbf{no training and nothing to converge}: the model \emph{is} the training
set. This is worth saying explicitly, because ``how do you know it has
converged?'' is a natural question and the answer is that there is no iteration
to converge. All the learning happened when we chose the projection.

We compare \emph{squared} distances and never take a square root, because
$\sqrt{\cdot}$ is increasing and so cannot change which neighbour is nearest.

\paragraph{The syntax.}
\begin{description}[leftmargin=*,style=nextline,nosep]
\item[\texttt{trainProj - testProj(:, i)}] Implicit expansion doing real work:
  one test column is subtracted from \emph{all} training columns at once,
  producing a whole matrix of differences without a loop.
\item[\texttt{.\textasciicircum 2}] \textbf{Element-wise} power. The dot is
  essential: \texttt{A\textasciicircum 2} would mean the matrix product $AA$.
\item[\texttt{sum(\ldots, 1)}] Sums down each column, collapsing the matrix of
  differences to one squared distance per training image.
\item[\texttt{[\textasciitilde, best] = min(d)}] We want \emph{where} the
  minimum is, not its value.
\end{description}

%% =====================================================================
\section{Part 3: the figures}

\begin{lstlisting}
figure('Visible', 'off');
imagesc(reshape(K.xbar, K.H, K.W));
colormap gray; axis image off;
title([K.name ': mean face']);
saveas(gcf, [figDir K.name '_meanface.png']);
close;
\end{lstlisting}

\paragraph{The maths.} Only one idea: a column of $d$ numbers is turned back
into an $H\times W$ picture. This is the exact inverse of \texttt{im(:)} from
Section~\ref{sec:loadfn}, and it works precisely because every image was
unrolled the same way. Note this means an eigenface or a Fisherface
\emph{can} be displayed as an image: they live in the same $\R^d$ as the faces,
even though they are directions rather than photographs, which is why they can
have negative entries.

\paragraph{The syntax.}
\begin{description}[leftmargin=*,style=nextline,nosep]
\item[\texttt{reshape(v,H,W)}] Refolds a vector into an $H\times W$ matrix,
  filling column by column --- the inverse of \texttt{v = im(:)}. The element
  count must match exactly.
\item[\texttt{figure('Visible','off')}] Creates a figure without popping a
  window, so the script runs headless.
\item[\texttt{imagesc(A)}] Draws a matrix as an image, \emph{scaling} the
  values to use the full colour range. (Plain \texttt{image} would clip
  instead, and eigenfaces have negative entries, so scaling is required.)
\item[\texttt{colormap gray}] Grey instead of the default colour map.
\item[\texttt{axis image off}] Square pixels, no axes or tick marks.
\item[\texttt{subplot(4,4,l)}] Splits the figure into a $4\times4$ grid and
  selects panel \texttt{l}, counted \emph{across rows first}.
\item[\texttt{sgtitle}] One title over a whole grid of subplots.
\item[\texttt{gcf}] ``Get current figure'' --- the handle \texttt{saveas} needs.
\item[\texttt{close}] Closes the figure. Without it a long loop would leave
  dozens of figures open and eventually exhaust memory.
\end{description}

%% =====================================================================
\section{Every construct in one table}
\label{sec:syntaxtable}

\begin{center}
\small
\begin{tabular}{llp{6.4cm}}
\toprule
construct & meaning & note \\
\midrule
\texttt{;} & suppress output & at end of a line \\
\texttt{\%} & comment & to end of line \\
\texttt{\ldots} & line continuation & statement continues below \\
\texttt{'} & transpose & \texttt{A'} is $A^{\!\top}$ \\
\texttt{*} vs \texttt{.*} & matrix vs element-wise product & the dot changes everything \\
\texttt{\textasciicircum} vs \texttt{.\textasciicircum} & matrix vs element-wise power & \texttt{A\textasciicircum2} is $AA$ \\
\texttt{A\textbackslash b} & solve $Ax=b$ & left solve, no inverse formed \\
\texttt{b/A} & solve $xA=b$ & right solve \\
\texttt{v(:)} & flatten to a column & column-major order \\
\texttt{A(:,j)} & column $j$ & colon $=$ ``all'' \\
\texttt{A(:,1:k)} & first $k$ columns & a range \\
\texttt{A(:,idx)} & reorder columns & \texttt{idx} a vector of positions \\
\texttt{A(:,y==j)} & logical indexing & keeps columns where the test is true \\
\texttt{[a, b]} & horizontal concat & \texttt{[a; b]} stacks vertically \\
\texttt{\{ \}} & cell array contents & \texttt{c\{i\}} gets the item out \\
\texttt{s.field} & struct field & \texttt{sets(1).name} \\
\texttt{[\textasciitilde, i] = f(\ldots)} & discard an output & keep only the second \\
\texttt{1:n} & the vector $1,\ldots,n$ & also used to drive \texttt{for} \\
\texttt{end} & closes \texttt{if}/\texttt{for}/\texttt{function} & one keyword for all \\
\midrule
\texttt{zeros(m,n)} & allocate & preallocate before loops \\
\texttt{size(A,k)} & length along dimension $k$ & \texttt{1} rows, \texttt{2} columns \\
\texttt{numel} & total element count & \\
\texttt{mean(A,1/2)} & average down columns / across & the dimension argument matters \\
\texttt{sum(A,1)} & sum down each column & \\
\texttt{norm(v)}, \texttt{norm(A,'fro')} & vector length, matrix size & \\
\texttt{min}, \texttt{max} & extremum & second output is \emph{where} \\
\texttt{sort(v,'descend')} & sort & second output is the permutation \\
\texttt{diag} & extract or build a diagonal & both, depending on input \\
\texttt{rank} & numerical rank & tolerance-based \\
\texttt{eig(A)} & eigenvalues/vectors & \textbf{ascending} order \\
\texttt{chol(A)} & Cholesky $A = R^{\!\top}R$ & \texttt{'lower'} gives $LL^{\!\top}$ \\
\texttt{randn}, \texttt{rng} & normal randoms, seed & \\
\texttt{erfc} & complementary error function & stands in for \texttt{normcdf} \\
\midrule
\texttt{dir(pattern)} & list files & struct array; sorts as \emph{text} \\
\texttt{imread}, \texttt{double} & read image, convert type & convert or integer maths bites \\
\texttt{reshape(v,H,W)} & vector back to image & inverse of \texttt{v(:)} \\
\texttt{imagesc}, \texttt{colormap} & draw a matrix & \texttt{imagesc} rescales \\
\texttt{subplot}, \texttt{sgtitle} & grid of panels, overall title & \\
\texttt{saveas}, \texttt{gcf}, \texttt{close} & write and tidy up figures & \\
\texttt{fprintf} & formatted printing & \texttt{\%d} integer, \texttt{\%f} float, \texttt{\%s} string, \texttt{\%\%} a literal percent \\
\texttt{error} & stop with a message & \\
\bottomrule
\end{tabular}
\end{center}
